% ============================================================
% f2012_guangdong_science_q06_views.tex
% 2012 广东卷（理科）第 6 题 三视图（TikZ 矢量重绘）
% 正视图、侧视图：下方矩形 6×5，上方等腰三角形底 6、腰 5（高 4）
% 俯视图：直径 6 的圆；还原为圆柱（r=3, h=5）与圆锥（r=3, 母线 5, h=4）
% ============================================================
\begin{tikzpicture}[
  x=1cm, y=1cm, line width=0.6pt, font=\footnotesize, >=Stealth,
  dim/.style={<->, line width=0.4pt}
]

  % ---------------- 正视图 ----------------
  \begin{scope}
    \draw (-1.5,0) rectangle (1.5,2.5);
    \draw (-1.5,2.5) -- (0,4.5) -- (1.5,2.5);
    % 底边 6
    \draw[dim] (-1.5,-0.42) -- (1.5,-0.42);
    \draw (-1.5,-0.34) -- (-1.5,-0.50);
    \draw (1.5,-0.34) -- (1.5,-0.50);
    \node[below] at (0,-0.42) {$6$};
    % 高 5
    \draw[dim] (1.88,0) -- (1.88,2.5);
    \draw (1.80,0) -- (1.96,0);
    \draw (1.80,2.5) -- (1.96,2.5);
    \node[right] at (1.90,1.25) {$5$};
    % 两腰 5
    \node[rotate=53.13] at (-1.20,3.30) {$5$};
    \node[rotate=-53.13] at (1.20,3.30) {$5$};
    \node at (0,-1.10) {正视图};
  \end{scope}

  % ---------------- 侧视图 ----------------
  \begin{scope}[xshift=5.4cm]
    \draw (-1.5,0) rectangle (1.5,2.5);
    \draw (-1.5,2.5) -- (0,4.5) -- (1.5,2.5);
    \draw[dim] (-1.5,-0.42) -- (1.5,-0.42);
    \draw (-1.5,-0.34) -- (-1.5,-0.50);
    \draw (1.5,-0.34) -- (1.5,-0.50);
    \node[below] at (0,-0.42) {$6$};
    \draw[dim] (1.88,0) -- (1.88,2.5);
    \draw (1.80,0) -- (1.96,0);
    \draw (1.80,2.5) -- (1.96,2.5);
    \node[right] at (1.90,1.25) {$5$};
    \node[rotate=53.13] at (-1.20,3.30) {$5$};
    \node[rotate=-53.13] at (1.20,3.30) {$5$};
    \node at (0,-1.10) {侧视图};
  \end{scope}

  % ---------------- 俯视图 ----------------
  \begin{scope}[yshift=-4.5cm]
    \draw (0,0) circle (1.5);
    \fill (0,0) circle (1.1pt);
    \node at (0,-1.95) {俯视图};
  \end{scope}

\end{tikzpicture}
